\documentclass[11pt]{article}

\usepackage[margin=0.9in]{geometry}
\usepackage{amsmath}
\usepackage{booktabs}
\usepackage{graphicx}
\usepackage{float}
\usepackage{hyperref}
\usepackage{array}
\usepackage[table]{xcolor}

\graphicspath{{../figures/}}
\sloppy

\newcommand{\mainfigure}[3]{%
\begin{figure}[H]
    \centering
    \includegraphics[width=\textwidth]{#1}
    \caption{#2}
    \label{#3}
\end{figure}
}

\title{FCFS Appointment Simulation: Metric-Focused Sensitivity Report}
\author{CUIMC Appointment Simulation}
\date{\today}

\begin{document}
\maketitle

\paragraph{Canonical source.} The canonical report source is
\texttt{docs/reports/metric\_analysis.qmd}. This TeX file is a generated or companion
surface and should not be edited as an independent report.

\section{Executive Summary}

This report is organized by metric. Each metric section asks the same questions:

\begin{enumerate}
    \item What does the metric mean?
    \item Which assumptions move it the most?
    \item What is the simplest plot to read?
    \item What do the two-class heatmaps add?
\end{enumerate}

\begin{table}[H]
\centering
\renewcommand{\arraystretch}{1.15}
\begin{tabular}{p{0.25\textwidth}p{0.29\textwidth}p{0.36\textwidth}}
\toprule
Metric & Main drivers & Main warning \\
\midrule
Average utilization & No-show risk, then cancellation and demand level & High utilization does not imply good access. \\
Overall served rate & Total demand, no-show risk, cancellations, balking & This is the main access metric; read it with utilization. \\
Mean offered wait & Total demand, balking tolerance, cancellations & Offered wait includes patients who later balk. \\
Balking rate & Balking step and balking threshold & Useful diagnostic; not a success metric by itself. \\
Class served-rate gap & Differences between class behavior parameters & Class labels matter only when class assumptions differ. \\
\bottomrule
\end{tabular}
\end{table}

The baseline treats the two classes symmetrically: Class 1 and Class 2 have the same arrival rate, cancellation probability, balking rule, no-show rule, and value. Symmetric heatmaps should therefore be read as absolute parameter effects, not class-label effects. Class advantage appears when the classes receive different behavior parameters.

\section{Baseline And Reading Order}

Baseline uses \texttt{configs/baseline.yaml}: $\lambda_1 = 50$, $\lambda_2 = 50$, 32 slots per day, a 14-day booking horizon, balking threshold 9, balking high probability 0.50, no-show threshold 6, no-show high probability 0.30, and cancellation probability 0.10 for both classes.

\begin{table}[H]
\centering
\begin{tabular}{lrrrrrr}
\toprule
Scenario & Utilization & Overall served & Accepted wait & Offered wait & Class gap & Delay gap \\
\midrule
Baseline & 0.839 & 0.269 & 8.35 & 9.30 & 0.001 & 0.003 \\
Scenario 2 & 1.000 & 0.395 & 4.29 & 5.06 & -0.008 & 0.005 \\
\bottomrule
\end{tabular}
\caption{Scenario values are rates or averages. Scenario 2 changes several assumptions at once, so it is a comparison point, not a one-parameter causal test.}
\end{table}

The plot order is intentional. Driver plots are simple one-axis checks: Class 1 changes on the x-axis while Class 2 stays fixed at its baseline value. Each driver plot shows the overall metric plus the Class 1 and Class 2 values when that class-level value is defined. Heatmaps are the more detailed view: the x-axis changes Class 1, the y-axis changes Class 2, and the dashed baseline line marks the current assumption when applicable.

\begin{table}[H]
\centering
\renewcommand{\arraystretch}{1.12}
\begin{tabular}{p{0.22\textwidth}p{0.68\textwidth}}
\toprule
Color family & Meaning \\
\midrule
Blue & Arrival pressure, arrival mix, and demand-load changes. \\
Purple & Balking step, balking threshold, and balking-rate diagnostics. \\
Green & No-show step and no-show threshold changes. \\
Red & Cancellation probability changes. \\
Line style and marker & Overall, Class 1, and Class 2 within the same driver family. \\
Gray dashed line & Baseline assumption in one-axis slices and heatmaps. \\
Driver-colored heatmaps & The varied assumption family. For class-gap heatmaps, lighter shades indicate Class 2 is higher and darker shades indicate Class 1 is higher. \\
\bottomrule
\end{tabular}
\caption{Color key used throughout the metric-focused figures.}
\end{table}

\section{Average Utilization}

\subsection*{Metric}

Average utilization measures completed visits per available slot:

\[
\texttt{average\_utilization} =
\frac{1}{D}\sum_d \frac{\text{served slots on day } d}{\text{slots per day}}.
\]

No-shows do not count as utilization because the slot did not become a completed visit. Overall utilization is aggregate. In the driver plot, the Class 1 and Class 2 lines decompose utilization into each class's share of all available slots.

\subsection*{What Impacts It}

No-show behavior is the clearest direct driver: booked slots disappear without becoming visits. Cancellation can also change utilization because it changes the appointment pool before service. Demand pressure is more subtle. Utilization can stay high across a wide demand range, even when the share of patients served collapses.

\mainfigure{metric_utilization_drivers.png}{Average utilization driver plots. Class 1 varies while Class 2 stays fixed at baseline. The class lines are each class's served slots divided by all available slots.}{fig:metric-util-drivers}

The simple view shows that utilization is not the same as access. Around baseline demand, the clinic is busy, but only a small share of arrivals are served. No-show assumptions move utilization more directly than demand once the system is already capacity constrained.

\mainfigure{no_show_step_interaction_heatmap.png}{No-show step heatmap. Left: utilization. Right: Class 1 served-rate gap.}{fig:metric-util-noshow-step}
\mainfigure{no_show_threshold_interaction_heatmap.png}{No-show threshold heatmap. Left: utilization. Right: Class 1 served-rate gap.}{fig:metric-util-noshow-threshold}

The heatmaps add the class-specific reading. When both classes have high no-show risk, utilization falls. When only one class has higher no-show risk, that class loses completed visits and the served-rate gap moves against it.

\section{Overall Served Rate}

\subsection*{Metric}

Overall served rate is the main access metric:

\[
\texttt{overall\_percent\_serviced} =
\frac{\sum_i V_i}{\sum_i A_i},
\]

where $A_i$ is arrivals and $V_i$ is served patients for class $i$.

\subsection*{What Impacts It}

The strongest aggregate driver is total arrival pressure. More arrivals compete for the same FCFS booking horizon and same daily capacity. No-shows and cancellations lower completed visits after booking. Balking lowers served rate because patients who reject long-delay offers leave the system.

\mainfigure{metric_access_drivers.png}{Served-rate driver plots. Class 1 varies while Class 2 stays fixed at baseline; each panel shows overall, Class 1, and Class 2 served rates.}{fig:metric-access-drivers}

The easiest reading is that served rate falls when demand exceeds capacity. At low demand nearly every arrival is served; near baseline the clinic remains highly utilized but access is already poor.

\mainfigure{fcfs_arrival_outcome_decomposition.png}{Final outcome decomposition under demand pressure. Each band is a share of arrivals.}{fig:metric-access-outcomes}

The outcome decomposition explains why served rate falls. At baseline, many arrivals are lost to balking or after-booking losses. Under heavier overload, \texttt{no\_offer} becomes a major failure mode because the FCFS horizon fills.

\mainfigure{arrival_mix_interaction_heatmap.png}{Arrival-rate and class-mix heatmap. Left: overall served rate. Right: offered wait.}{fig:metric-access-arrival-heatmap}
\mainfigure{cancellation_probability_interaction_heatmap.png}{Cancellation heatmap. Left: overall served rate. Right: Class 1 served-rate gap.}{fig:metric-access-cancel-heatmap}
\mainfigure{balking_step_interaction_heatmap.png}{Balking step heatmap. Left: overall served rate. Right: Class 1 served-rate gap.}{fig:metric-access-balk-step-heatmap}

The heatmaps show two different mechanisms. Arrival pressure mostly changes the aggregate served rate. Cancellation and balking differences can change both aggregate access and which class receives more completed visits.

\section{Mean Offered Booking Delay}

\subsection*{Metric}

Mean offered booking delay is the patient-facing wait metric:

\[
\texttt{mean\_offered\_booking\_delay} =
\frac{\sum_i \text{total offered delay}_i}{\sum_i O_i},
\]

where $O_i = B_i + K_i$ is offered patients: accepted bookings plus balked patients. Patients with \texttt{no\_offer} are excluded because they never receive a slot offer.

\subsection*{What Impacts It}

Demand pressure raises offered wait because patients are pushed farther into the booking horizon. Balking behavior changes wait by changing whether long-delay offers remain in the accepted appointment pool. Cancellations can shorten offered wait by freeing future capacity, but they do so by losing already-booked patients.

\mainfigure{metric_wait_drivers.png}{Mean offered-wait driver plots. Class 1 varies while Class 2 stays fixed at baseline; each panel shows overall, Class 1, and Class 2 offered waits.}{fig:metric-wait-drivers}

The simple view shows that offered wait is most intuitive under demand pressure: more arrivals push offers farther out. Balking and cancellation need more care because shorter waits can be caused by patients leaving the system, not by better access.

\mainfigure{arrival_rate_slice_wait.png}{Offered-wait slice as total arrivals change and class mix stays 50/50.}{fig:metric-wait-arrival-slice}
\mainfigure{balking_threshold_interaction_heatmap.png}{Balking threshold heatmap. Left: offered wait. Right: Class 1 served-rate gap.}{fig:metric-wait-balk-threshold}

The balking-threshold heatmap adds the class interpretation. The class with the higher tolerance threshold accepts longer waits and is more likely to stay in the system. That can increase its served rate while also keeping the system more congested.

\section{Balking Rate}

\subsection*{Metric}

Balking rate is a diagnostic for rejected offers:

\[
\texttt{balking\_rate} =
\frac{\sum_i K_i}{\sum_i O_i}.
\]

It is measured among patients who received an offer. No-offer patients are excluded because they never had the opportunity to accept or reject a delay.

\subsection*{What Impacts It}

Balking rate is controlled directly by the balking step and threshold. A higher step makes high-delay offers more likely to be rejected. A lower threshold makes that high rejection probability start earlier.

\mainfigure{metric_balking_rate_drivers.png}{Balking-rate driver plots. Class 1 varies while Class 2 stays fixed at baseline; each panel shows overall, Class 1, and Class 2 balking rates.}{fig:metric-balk-rate-drivers}

Balking rate helps explain why a lower wait is not always an improvement. If wait falls because many long-delay offers are rejected, access can worsen.

\mainfigure{balking_step_balking_rate_heatmap.png}{Balking-rate heatmap under balking-step changes. Left: overall balking rate. Right: Class 1 minus Class 2 balking-rate gap.}{fig:metric-balk-rate-step-heatmap}
\mainfigure{balking_threshold_balking_rate_heatmap.png}{Balking-rate heatmap under balking-threshold changes. Left: overall balking rate. Right: Class 1 minus Class 2 balking-rate gap.}{fig:metric-balk-rate-threshold-heatmap}

The heatmaps separate absolute balking from class differences. When both classes have high balking, overall balking rises. When one class has higher balking than the other, the more patient class tends to receive more service.

\section{Class Served-Rate Gap}

\subsection*{Metric}

Class served-rate gap is the main class-benefit metric:

\[
\Delta_{\text{access}} =
\texttt{percent\_serviced}_1 - \texttt{percent\_serviced}_2.
\]

Positive values mean Class 1 is served more often. Negative values mean Class 2 is served more often.

\subsection*{What Impacts It}

Class gap is driven by differences between class assumptions. In the current FCFS baseline, class labels are otherwise exchangeable. Cancellation gaps are the largest observed class-access driver, followed by balking threshold gaps, balking step gaps, and no-show gaps.

\mainfigure{metric_class_gap_drivers.png}{Served-rate values behind class gaps. Class 1 changes while Class 2 stays fixed at baseline. The class gap is the vertical distance between the Class 1 and Class 2 lines.}{fig:metric-class-gap-drivers}

The simple class-gap plots show direction through the distance between class lines. Higher Class 1 cancellation probability, higher Class 1 balking step, and higher Class 1 no-show step move the Class 1 line below the Class 2 line. A higher Class 1 balking threshold helps Class 1 because it tolerates longer offered waits.

\mainfigure{cancellation_probability_interaction_heatmap.png}{Cancellation interaction heatmap. The right panel shows the served-rate gap created by cancellation differences.}{fig:metric-class-gap-cancel}
\mainfigure{no_show_step_interaction_heatmap.png}{No-show step interaction heatmap. The right panel shows the served-rate gap created by no-show differences.}{fig:metric-class-gap-noshow}
\mainfigure{balking_threshold_interaction_heatmap.png}{Balking threshold interaction heatmap. The right panel shows the served-rate gap created by tolerance differences.}{fig:metric-class-gap-balk-threshold}

The heatmaps show symmetry around the baseline because the baseline classes are identical. Moving to the upper-left or lower-right corners creates opposite class advantages.

\section{Regression Screen}

The line plots and heatmaps isolate one or two assumptions at a time. The regression screen checks the same conclusions across 240 random FCFS parameter settings, with two simulation seeds averaged per setting. Inputs include total demand, class mix, balking steps and thresholds, no-show steps and thresholds, and cancellation probabilities.

\begin{table}[H]
\centering
\begin{tabular}{llr}
\toprule
Target metric & Most important feature & Standardized coefficient \\
\midrule
Utilization & Average no-show threshold & 0.512 \\
Utilization & Average no-show step & -0.423 \\
Utilization & Average cancellation probability & 0.320 \\
Overall served rate & Total arrival rate & -0.774 \\
Overall served rate & Average no-show threshold & 0.251 \\
Offered wait & Total arrival rate & 0.576 \\
Offered wait & Average cancellation probability & -0.441 \\
Class gap & Cancellation probability gap & -0.452 \\
Class gap & Balking threshold gap & 0.429 \\
Class gap & Balking step gap & -0.349 \\
\bottomrule
\end{tabular}
\caption{Top regression drivers. Gap means Class 1 value minus Class 2 value.}
\end{table}

\mainfigure{regression_standardized_coefficients.png}{Full standardized regression coefficient plot from randomized FCFS simulations.}{fig:metric-regression-full}

The regression screen supports the metric-by-metric reading: average no-show behavior explains utilization, total arrival rate explains access and wait, and class parameter gaps explain class advantage.

\section{Final Takeaways}

\begin{enumerate}
    \item Use \texttt{overall\_percent\_serviced} as the main access metric and \texttt{average\_utilization} as the capacity-use metric.
    \item Do not interpret high utilization as good performance by itself. In the baseline, utilization is about 0.839 while only about 0.269 of arrivals complete service.
    \item Use \texttt{mean\_offered\_booking\_delay} for patient-facing wait because it includes patients who balk after seeing a long-delay offer.
    \item Use balking rate as a diagnostic, not as a final success metric.
    \item For class benefit, look at behavior gaps. Cancellation differences are the strongest class-gap driver in this run, with balking and no-show differences also important.
\end{enumerate}

\appendix
\section{Appendix: Exact Metric Definitions}

Let $i$ be patient class, $D$ be measured days, $S$ be slots per day, $A_i$ be arrivals, $B_i$ be booked patients, $K_i$ be balked patients, $C_i$ be canceled patients, $H_i$ be no-shows, $V_i$ be served patients, and $O_i = B_i + K_i$ be offered patients.

\begin{align*}
\texttt{percent\_serviced}_i &= \frac{V_i}{A_i}, \\
\texttt{overall\_percent\_serviced} &= \frac{\sum_i V_i}{\sum_i A_i}, \\
\texttt{mean\_offered\_booking\_delay} &= \frac{\sum_i \text{total offered delay}_i}{\sum_i O_i}, \\
\texttt{mean\_accepted\_booking\_delay} &= \frac{\sum_i \text{total accepted delay}_i}{\sum_i B_i}, \\
\texttt{average\_utilization} &= \frac{1}{D}\sum_d \frac{\text{served slots}_d}{S}, \\
\texttt{balking\_rate} &= \frac{\sum_i K_i}{\sum_i O_i}, \\
\Delta_{\text{access}} &= \texttt{percent\_serviced}_1 - \texttt{percent\_serviced}_2.
\end{align*}

\section{Appendix: Dense Reference Plots}

These figures are retained for traceability when the full diagnostic view is useful.

\mainfigure{scenario_metric_comparison.png}{Scenario-level metric comparison.}{fig:appendix-scenario}
\mainfigure{fcfs_capacity_stress_curves.png}{Full FCFS stress diagnostic plot.}{fig:appendix-fcfs-full}
\mainfigure{balking_threshold_jump_heatmaps.png}{Balking threshold and jump level when both classes share the same values.}{fig:appendix-balk-common}
\mainfigure{no_show_threshold_jump_heatmaps.png}{No-show threshold and jump level when both classes share the same values.}{fig:appendix-noshow-common}

\end{document}
